

\part{Definable construction sequences}\label{part:definable-construction-sequences}

Recall the following definitions from \Cref{sec:skel-part-defseq}.
\defpartialdef*

\defskellocdef*

In this part, we establish the following main result.
\thmskeldefseq*

We will use the following variation on the definition of merge-width, originating in~\cite{drabik2026variants}.

% lax begin Lax185574.Chains
\begin{definition}\label{def:chain-width}
  A \emph{chain} in a graph $G$ is a decreasing sequence $V(G)=S_1\supseteq\dots\supseteq S_\ell=\emptyset$ of vertex sets. Its \emph{radius-$r$ width}, for $r\in\N$, is defined as
$$\max_{2\le i\le \ell} \max_{v\in V(G)}\left|\ATypes^G(\Ball_r^{S_{i}}(v)/S_{i-1})\right|=\max_{2\le i\le \ell} \max_{v\in V(G)}\left|\Ball_r^{S_{i}}(v)/\PP_{S_{i-1}}\right|,$$
where $\Ball_r^S(v)$ is defined above via $S$-flips and $\PP_{S}$ is the partition into atomic types over $S$.
\end{definition}

% lax end

The offset between $S_{i-1}$ and $S_i$ in the definition of the radius-$r$ width ensures that the definition is robust. For instance, a chain can be completed into one in which two consecutive sets differ by one element,
by removing the elements of $S_{i-1}$ one by one in any order, until $S_i$ is reached; this transformation
preserves the width parameters.

Chains can be converted to construction sequences and vice versa, with a loss in the radius-$r$ width.
Intuitively, a chain
gives rise to a sequence of ever coarser partitions $\PP_{S_1}\preceq \PP_{S_2}\preceq\dots \preceq \PP_{S_\ell}$, which can then be converted
into a sequence of partial graphs
(see~\cite{drabik2026variants} for details).

% lax begin Lax185574Proofs.ComputingChains.exists_program_chain
We use the following key result from~\cite{md-abmw}. It says that in monadically dependent classes, chains of constant length and small width can be computed efficiently.
% lax end

% lax begin Lax185574.ComputingChains
\begin{theorem}[{\cite[Theorem~4.4]{md-abmw}}]\label{thm:monadic-dependence-implies-definable-mw}
  There is an algorithm which, given
  an $n$-vertex graph $G$, radius $r\in\N$, and rational $\eps>0$, computes a chain of $G$ of length $O(r/\eps)$ and radius-$r$ width $O_{\CC,r,\eps}(n^\eps)$,
  and runs in time  $\widehat O_{\CC,r}(n^{3})$,
  where $\CC$ is any monadically dependent graph class containing $G$.
\end{theorem}
% lax end

To prove \Cref{thm:skel-def-seq}, we turn these chains into construction sequences using the following main technical contribution of this part.

% lax begin Lax185574.ChainConstructionSequences
\begin{restatable}{lemma}{lembetterdefinablemwfromsets}\label{lem:better-definable-mw-from-sets}
  For every $n$-vertex graph $G$ and its chain of length $\ell$ there is
  a construction sequence $G_1,\ldots,G_\lambda$ of $G$ such that the following properties hold for all $r\ge 1$.
  Let $k_r$ denote the radius-$(5r+1)$ width of the chain and $w_r \coloneqq k_r\cdot \pi_G(k_r)$.
  \begin{enumerate}
    \item The resulting construction sequence has radius-$r$ width at most $w_r$ and length $\lambda\le O(nw_1)$.
    \item For every $t\in[\lambda-1]$ and $v\in V(G)$, the induced partial graph
    \[
        G_t[\Ball^{G_{t+1}}_r(v)]
    \]
    is definable in some expansion of $G$
    with $2\ell+1$ unary predicates, using formulas of quantifier rank~$7$ with $w_r$ parameters.
    \item The construction sequence can be computed in time $O(\ell n^3)$, given $G$ and its chain.
  \end{enumerate}
\end{restatable}
% lax end

% lax begin Lax185574Proofs.DefinableSequences.exists_program_definableSequence
\begin{proof}[Proof of \Cref{thm:skel-def-seq}]
  Let $\CC$ be a monadically dependent graph class containing $G$, and let $d$ be an upper bound on the VC-dimension of graphs in $\CC$, which exists by \Cref{obs:md-vc}.
  We apply \Cref{thm:monadic-dependence-implies-definable-mw} with radius $5r+1$ and $\eps$.
  This computes a chain of length $\ell\le O((5r+1)/\eps)=O(r/\eps)$ in time $\widehat O_{\CC,r}(n^3)$.
  We then apply \Cref{lem:better-definable-mw-from-sets} to $G$ and this chain, which takes time $O(\ell n^3)$.
  So the total running time is
  \[
    \widehat O_{\CC,r}(n^3)+O(\ell n^3)
    =\widehat O_{\CC,r,\eps}(n^3).
  \]

  The chain has radius-$(5r+1)$ width $k_r\le O_{\CC,r,\eps}(n^\eps)$, and $\pi_G(k_r)\le(k_r+1)^d$
  by the Sauer--Shelah lemma (\Cref{lem:ss}), so
  \[
    w_r=k_r\cdot\pi_G(k_r)\le(k_r+1)^{d+1}\le O_{\CC,r,\eps}(n^{(d+1)\eps}).
  \]
  We set $\delta\coloneqq d+1$.
  The sequence is radius-$r$ definable with $2\ell+1\le O(r/\eps)$ unary predicates, $w_r$ parameters, and rank $7$.
  As $k_1\le k_r$, it has length $O(nw_1)\le O(nw_r)\le O_{\CC,r,\eps}(n^{1+\delta\eps})$.
\end{proof}
% lax end

\paragraph{Proof strategy for \texorpdfstring{\Cref{lem:better-definable-mw-from-sets}}{the technical lemma}}\label{sec:defseq-lemma-strategy}

The proof of \Cref{lem:better-definable-mw-from-sets}
occupies the remainder of \Cref{part:definable-construction-sequences}.
It combines constructions from~\cite{drabik2026variants} and~\cite{mw-stoc}, together with a careful analysis of definability.
We revisit several variants of merge-width studied in those papers.
In particular, we use the following variation of the definition of merge-width.

% lax begin Lax185574.MergeSequences
\begin{definition}\label{def:merge-sequence}
A \emph{merge sequence} of a graph $G$ is a sequence
$(\PP_1,R_1),\ldots,(\PP_\ell,R_\ell)$ such that:
\begin{enumerate}
  \item $\PP_1$ is the partition into
singletons and $\PP_\ell=\{V(G)\}$,
\item $\PP_1\preceq \dots\preceq \PP_\ell$,
\item $R_1\subseteq \dots\subseteq R_\ell\subseteq {V(G)\choose 2}$,
\item for every $i\in[\ell]$ and $X,Y\in \PP_i$,
the set $XY-R_i$ is contained in either $E(G)$ or in $E(\overline G)$,
 that is, the pair $X,Y$ is homogeneous once pairs in $R_i$ are ignored.
\end{enumerate}

  The \emph{radius-$r$ width} of a merge sequence as above is
\begin{equation}\label{eq:merge-seq-width}
  \max_{2\le i\le \ell}\max_{v\in V(G)}
  \left|\Ball^{R_i}_r(v)/\PP_{i-1}\right|.
\end{equation}
\end{definition}
% lax end
\noindent
Here the offset in the indices of $R_i$ and $\PP_{i-1}$ is intentional.
Every construction sequence naturally induces a merge sequence of the same radius-$r$ width for all $r\in\N$, and vice versa (see~\cite{mw-arxiv,mw-stoc} for details).

\section{Merge-width with definable stages}\label{sec:better-definable-mw}
\label{sec:definable-mw}

In this section, we prove the following lemma.

% lax begin Lax185574.BetterMergeSequences
\begin{restatable}{lemma}{lembettermw}\label{lem:better-mw}
For every graph $G$ and chain of length $\ell$ there is a merge sequence $(\QQ_1, R_1), \ldots, (\QQ_\ell, R_\ell)$ with the following properties.
    \begin{enumerate}
        \item For all $r\ge 1$,  the radius-$r$ width of the merge sequence is bounded by $k\cdot \pi_G(k)$, where $k$ is the radius-$(5r+1)$ width of the given chain.
        \item  For each $i\in [\ell]$, the relations $\equiv_{\QQ_i}$ and $R_i$ are definable in some expansion $G^i$ of $G$ with two unary predicates, using two fixed first-order formulas $\phi_\equiv(x,y)$ and $\phi_R(x,y)$ of quantifier rank at most~$5$.
        \item The merge sequence can be computed in time $O(\ell n^3)$, given $G$ and its chain.
    \end{enumerate}
\end{restatable}
% lax end

\subsection{Preliminaries}

The proof of \Cref{lem:better-mw} uses ideas developed in~\cite{drabik2026variants}. We first recall the required notions and lemmas from that paper.

% lax begin Lax185574.ComplementDiameter
\begin{lemma}[folklore]\label{lem:diamComplement}
    For every graph $G$, at least one of $G$ and $\overline{G}$ has diameter at most $3$.
\end{lemma}
% lax end
% lax begin Lax185574Proofs.ComplementDiameter.diameter_le_three_or_compl
\begin{proof}
Suppose $G$ has diameter greater than $2$. Then there are non-adjacent vertices $u,v$ with no common neighbor in $G$.
In $\overline G$, $uv$ is an edge and all vertices are neighbors of either $u$ or $v$.
Thus $\overline G$ has diameter at most $3$.
\end{proof}
% lax end

% lax begin Lax185574.BipartiteGraphs
If $B$
is a bipartite graph with parts $X,Y$ and edge set $E$,
its \emph{bipartite complement}, denoted $\overline B$,
is the bipartite graph with parts $X,Y$ and edge set $XY\setminus E$.
% lax end

% lax begin Lax185574.BipartiteFlip
\begin{lemma}[{\cite[Lemma~A.3]{drabik2026variants}}]\label{lem:bipartite-informal-lemma}
  Let $B = (X,Y,E)$ be a bipartite graph with bipartite complement~$\overline{B}$, and pick arbitrary vertices $x \in X$, $y \in Y$.
  Partition~$X$ according to adjacency to $y$,
  i.e.\ $X = X_+ \uplus X_-$ with $X_+ = N(y)$ and $X_- = X \setminus N(y)$,
  and similarly $Y = Y_+ \uplus Y_-$ according to~$x$.
  Then there is a bipartite graph $B' = (X,Y,E')$ obtained from $B$ by some $\{X_+,X_-,Y_+,Y_-\}$-flip satisfying:
  \[ \text{for all $uv \in E'$,} \quad \dist^B(u,v) \le 5 \quad \text{and} \quad \dist^{\overline{B}}(u,v) \le 5. \]
\end{lemma}
% lax end
% lax begin Lax185574Proofs.BipartiteFlip.exists_flip_close
Note that~\cite[Lemma~A.3]{drabik2026variants} states the above lemma with $6$ instead of $5$; however, the two statements are equivalent, since
if $uv\in E'$ then $u$ and $v$ lie in opposite parts of the bipartite graphs $B'$, $B$, and $\overline{B}$, and thus $\dist(u,v)\le 6$ is equivalent to $\dist(u,v)\le 5$ for parity reasons.
% lax end

% lax begin Lax185574.TypeSplitting
\begin{definition}
    Fix a graph $G$ and a partition $\PP$ of $V(G)$. We define $\PP \wedge T$ to be the refinement of~$\PP$ where each part $P \in \PP$ is split into atomic types over $T \setminus P$
    (i.e.\ into sets of the form $P\cap Q$, where $Q$ is the set of all vertices of $G$ with a given neighborhood in $T-P$).
\end{definition}
% lax end

\noindent In this definition, we choose not to divide a part $P$ into types over $T \cap P$. This is for technical convenience only, to ensure that $\PP \wedge T = \PP$ whenever $\PP$ is the partition with only one part.

% lax begin Lax185574.SplitBalls
% lax begin Lax185574Proofs.SplitBalls.partCount_split_le
\begin{lemma}[{\cite[Corollary~5.4]{drabik2026variants}}]\label{lem:transversalRefinement}
	Let~$G $ be a graph, $\PP\preceq\QQ$ partitions of~$V(G)$, and~$T$ a transversal of~$\PP$.
  Suppose $|\Ball^\QQ_{r+1}(v)/\PP| \le k$ for every $v\in V(G)$.
  Then for every $v\in V(G)$, we have
  \[
    \left|\Ball^{\QQ}_r(v)/(\PP \wedge T)\right|\le k \cdot \pi_G(k).
  \]
\end{lemma}
% lax end
% lax end

\subsection{Constructing the merge sequence}

% lax begin Lax185574.BipartiteGraphs
Recall that for a graph
$G$ and disjoint vertex sets $X,Y\subseteq V(G)$,
by $G[X,Y]$ we denote the bipartite graph with parts $X$ and $Y$, and edges connecting pairs of vertices which are adjacent in $G$.
Furthermore, $G[X,X]$
denotes the bipartite graph whose parts are two copies of $X$,
and edges connect two vertices in opposite parts if the corresponding vertices in $G$ are either adjacent or equal.
We call the edges of $G[X,X]$ connecting the two copies of each vertex in $X$ \emph{diagonal edges}.
Thus,
${\overline G[X,X]}$ is the bipartite complement  $\overline {G[X,X]}$ of $G[X,X]$ extended additionally with the diagonal edges.
% lax end

% lax begin Lax185574.PairRelation
\begin{definition}
    Let $G$ be a graph and $\PP$ a partition of $V = V(G)$.
    Define $R = R(G,\PP) \subseteq \binom{V}{2}$ so that for all parts $X,Y \in \PP$ (possibly $X = Y$) and pairs $uv$ with $u \in X$, $v \in Y$, and $u \ne v$, we have
    $uv \in R$ if and only if
    \[
    \dist^{G[X,Y]}(u,v) \le 5 \quad \text{and} \quad \dist^{\overline{G}[X,Y]}(u,v) \le 5. \]
\end{definition}
% lax end

We first describe the construction of the merge sequence that underlies the proof of \Cref{lem:better-mw}.
Consider a graph $G$ and its chain
    $V(G)=S_1\supseteq\cdots\supseteq S_\ell=\emptyset$.
    Let $\PP_i$ be the partition into atomic types over $S_i$.
    Fix an arbitrary linear order $<$ on $V(G)$
    giving rise to representative sets $T_i \coloneqq \{\min_< P \,:\, P \in \PP_i\}$.
    Then $T_1\supseteq T_2\supseteq \dots \supseteq T_\ell$.
    For $i=1,\ldots,\ell$ define $$R_i \coloneqq R(G,\PP_i)\qquad\text{ and }\qquad\QQ_i \coloneqq \PP_i \wedge T_i.$$
    The following lemmas establish that
    $$(\QQ_1,R_1),\ldots,(\QQ_\ell,R_\ell)$$ is a merge sequence satisfying the conditions stated in \Cref{lem:better-mw}.

% lax begin Lax185574.PairRelationDefinable
\begin{lemma}\label{lem:def-R}
    Suppose $G$ is a graph and $\PP$ is the partition of $V(G)$ into atomic types over some set $S \subseteq V(G)$. Then the relation $R \coloneqq R(G,\PP)$ is definable in $G$ expanded with one unary predicate marking~$S$, using a fixed first-order formula $\phi_R(x,y)$ of quantifier rank at most~$5$.
\end{lemma}
% lax end

% lax begin Lax185574Proofs.PairRelationDefinable.exists_formula_pairRelation
\begin{proof}
    The following formula $\theta_S(x,y)$, using a unary predicate marking $S$ and with one quantifier,  defines the equivalence relation $\equiv_\PP$:
    \begin{align}\label{eq:definableEquiv}
        \theta_S(x,y) \coloneqq
        x = y \lor \bigl(\neg S(x) \land \neg S(y) \land
        \forall z\,(S(z) \to (E(x,z) \leftrightarrow E(y,z)))\bigr).
    \end{align}

  The following formula $\eta_S(x,y;a,b)$ says that $x$ and $y$ are in opposite parts of the bipartite graph $G[X,Y]$, where $X$ and $Y$ are the parts of $\PP$ containing $a$ and $b$:
\begin{align*}
  \eta_S(x,y;a,b)\coloneqq{}&
  (\theta_S(x,a)\land\theta_S(y,b))
  \lor(\theta_S(y,a)\land\theta_S(x,b)).
\end{align*}
The following formulas $e_S^+(x,y;a,b)$ and $e_S^-(x,y;a,b)$
 say that $x$ and $y$ are adjacent in the bipartite graph $G[X,Y]$ or $\overline G[X,Y]$, respectively, where $X$ and $Y$ are the parts of $\PP$ containing $a$ and $b$:
\begin{align*}
  e_S^+(x,y;a,b)&\coloneqq \eta_S(x,y;a,b)\land(x=y\lor E(x,y)),\\
  e_S^-(x,y;a,b)&\coloneqq \eta_S(x,y;a,b)\land(x=y\lor \neg E(x,y)).
\end{align*}
For $\sigma\in\{+,-\}$, define
\[
  \rho_S^\sigma(x,y)\coloneqq
  \bigvee_{p\in\{1,3,5\}}
  \exists z_1\cdots\exists z_{p-1}\,
  \bigwedge_{i=0}^{p-1}e_S^\sigma(z_i,z_{i+1};x,y),
\]
where $z_0\coloneqq x$ and $z_p\coloneqq y$.
This formula expresses the existence of a walk of length $1$, $3$, or $5$
between opposite copies of $x$ and $y$ in $G[X,Y]$ for $\sigma=+$, and
in $\overline G[X,Y]$ for $\sigma=-$, where $X\coloneqq \PP(x)$ and $Y\coloneqq \PP(y)$.
Each step switches sides, also when $X=Y$; an equality step uses a diagonal edge.
Since every path between opposite sides of a bipartite graph has odd length,
these formulas express precisely the two distance bounds defining $R$.
Consequently,
\[
  \phi_R(x,y)\coloneqq
  x\ne y\land\rho_S^+(x,y)\land\rho_S^-(x,y)
\]
defines $R(G,\PP)$.
Each $e_S^\sigma$ has quantifier rank~$1$, and each disjunct in
$\rho_S^\sigma$ quantifies at most four intermediate vertices.
Thus $\phi_R$ is a fixed first-order formula of quantifier rank at most~$5$.
\end{proof}
% lax end

% lax begin Lax185574.PairRelationMonotone
\begin{lemma}\label{lem:monotoneR}
    Fix a graph $G$ and partitions $\PP$ and $\PP'$ of $V(G)$ such that $\PP$ refines $\PP'$. Then $$R(G,\PP)\subseteq R(G, \PP').$$
\end{lemma}
% lax end
% lax begin Lax185574Proofs.PairRelationMonotone.pairRelation_mono
\begin{proof}
    Choose any $uv \in R(G,\PP)$. Let
    \[
        X \coloneqq \PP(u), \quad
        X' \coloneqq \PP'(u), \quad
        Y \coloneqq \PP(v), \quad
        Y' \coloneqq \PP'(v)
    \]
    be the parts of $\PP$ and $\PP'$ that contain $u$ and $v$.
    As $\PP$ refines $\PP'$, we have $X \subseteq X'$ and $Y \subseteq Y'$.
    Hence $G[X,Y]$ and $\overline G[X,Y]$ are induced subgraphs of $G[X',Y']$ and $\overline G[X',Y']$, respectively.
    Thus, by the definition of $R$,
    \begin{align*}
        \dist^{G[X',Y']}(u,v) &\le \dist^{G[X,Y]}(u,v) \le 5,\\
        \dist^{\overline G[X',Y']}(u,v) &\le \dist^{\overline G[X,Y]}(u,v) \le 5.
    \end{align*}
    This proves $uv \in R(G, \PP')$, as required.
\end{proof}
% lax end

% lax begin Lax185574.PairRelationBalls
\begin{lemma}\label{lem:ballContainment}
    Fix a graph $G$ and a partition $\PP$ of $V(G)$ and let $R \coloneqq R(G,\PP)$. For every radius $r \in \N$ and vertex $v \in V(G)$ we have
    \[
        \Ball^R_r(v) \subseteq \Ball^\PP_{5r}(v).
    \]
\end{lemma}
% lax end

% lax begin Lax185574Proofs.PairRelationBalls.graphBall_pairRelation_subset
\begin{proof}
    It suffices to show that for every $uv \in R$ we have $\dist_\PP(u,v) \leq 5$.
    Let $X,Y\in\PP$ be the parts containing $u$ and $v$, respectively. Fix an arbitrary $\PP$-flip $H$ of $G$.
    Then the bipartite graph $H[X,Y]$ is either $G[X,Y]$ or $\overline G[X,Y]$.
    By the definition of $R$, this graph contains a $u$--$v$ path of length at most $5$.
    If $X \neq Y$, then this path is also a path in $H$.
    If $X = Y$, we can project this path to a length ${\leq}5$ path in $H$ by contracting diagonal edges. Hence, in both cases,
    $\dist^H(u,v) \leq 5$, which finishes the proof.
\end{proof}
% lax end

% lax begin Lax185574.HomogeneousSplitting
\begin{lemma}\label{lem:default}
    Fix a graph $G$ and a partition $\PP$ of $V(G)$ and let $T$ be a transversal of $\PP$ and $R \coloneqq R(G,\PP)$.
    Then $\QQ \coloneqq \PP \wedge T$ is \emph{homogeneous modulo $R$}, i.e., for all $X,Y \in \QQ$ either
    \[
        XY \setminus R \subseteq E(G)
        \quad
        \text{or}
        \quad
        XY \setminus R \subseteq E(\overline{G}).
    \]
\end{lemma}
% lax end

% lax begin Lax185574Proofs.HomogeneousSplitting.homogeneousModulo_split
\begin{proof}
    We construct a $\QQ$-flip $H$ of $G$ with $E(H)\subseteq R$. This will yield the conclusion.
    Indeed, let $X,Y\in\QQ$. Since $H$ is a $\QQ$-flip of $G$, either
    \begin{center}
    $E(H)\cap XY \,=\, E(G)\cap XY$ \quad or \quad  $E(H)\cap XY \,=\, E(\overline G)\cap XY$.
    \end{center}
    Since $E(H)\subseteq R$, no pair in $XY\setminus R$ is an edge of $H$.
    Hence $XY\setminus R\subseteq E(\overline G)$ in the first case, and
    $XY\setminus R\subseteq E(G)$ in the second.
    It remains to define $H$. We separately describe how to flip a part $X$ with itself, and how to flip two distinct parts $X$ and $Y$ with each other.

    \begin{itemize}
        \item     For every part $X\in\PP$, by \Cref{lem:diamComplement} we have that
    either $G[X]$ or $\overline{G}[X]$ has diameter at most~$3$.
    If $G[X]$ has diameter at most $3$, set
 $H[X]\coloneqq \overline G[X]$, otherwise set $H[X]\coloneqq G[X]$.

 Suppose now that $uv\in E(H[X])$ for some $u,v\in X$.
 We consider the case when $G[X]$ has diameter at most $3$ and
  $H[X] = \overline G[X]$, the other case being symmetric.
Then $\dist^{G[X]}(u,v)\le 3$ and $\dist^{\overline G[X]}(u,v)= 1$.
The path of length at most $3$ in $G[X]$ lifts to a path of length at most $3$ between opposite copies in $G[X,X]$ after adding at most one diagonal step to make its length odd.
Similarly, the edge $uv$ in $\overline G[X]$ lifts to a path of length $1$ between opposite copies in $\overline G[X,X]$. This shows that $uv\in R$, and thus $E(H[X])\subseteq R$.

\item
For every pair of distinct parts $X,Y\in\PP$, let $x,y$ be their representatives in $T$ and apply \Cref{lem:bipartite-informal-lemma} to $G[X,Y]$.
    Use the resulting bipartite graph as $H[X,Y]$. The statement of the lemma implies that $E(H[X,Y])\subseteq R$.

    The partition $\QQ$ restricted to $X\cup Y$, denote it by $\QQ[X,Y]$, refines the neighborhood partitions of $X$ over $y$ and of $Y$ over $x$, so $H[X,Y]$ is a $\QQ[X,Y]$-flip of $G[X,Y]$.
    Altogether, $H$ is a
    $\QQ$-flip of $G$ with $E(H)\subseteq R$.\qedhere
    \end{itemize}
\end{proof}
% lax end

We now prove \Cref{lem:better-mw}.

\lembettermw*

% lax begin Lax185574Proofs.BetterMergeSequences.exists_mergeSequence_of_chain
\begin{proof}
    Let
    $V(G)=S_1\supseteq\cdots\supseteq S_\ell=\emptyset$ be the given chain.
    Let $\PP_i$ be the partition into atomic types over $S_i$.
    Fix an arbitrary linear order $<$ on $V(G)$
    giving rise to representative sets $T_i \coloneqq \{\min_< P \,:\, P \in \PP_i\}$.
    Then $T_1\supseteq T_2\supseteq \dots \supseteq T_\ell$.
    Define $$R_i \coloneqq R(G,\PP_i)\qquad\text{ and }\qquad\QQ_i \coloneqq \PP_i \wedge T_i.$$

    We claim that
    \[
        (\QQ_1, R_1), \ldots, (\QQ_\ell, R_\ell)
    \]
    is a merge sequence.  Indeed:
    \begin{enumerate}
        \item Since $\PP_1$ is the partition into singletons, so is
        $\QQ_1$.
        It follows from $\PP_\ell=\{V(G)\}$ and the definition of $\PP_\ell \wedge T_\ell$ that also
        $\QQ_\ell=\{V(G)\}$.

% lax begin Lax185574.SplittingMonotone
% lax begin Lax185574Proofs.SplittingMonotone.split_minimalVertices_mono
        \item For $i<j$ we have that $T_i\supseteq T_j$ and $\PP_i\preceq \PP_j$, which implies
that  $\PP_i\wedge T_i\preceq \PP_j\wedge T_j$. Hence,

    \[
        \QQ_1 \preceq \QQ_2 \preceq \dots \preceq \QQ_\ell
    \]
    is a sequence of ever coarser partitions.
% lax end
% lax end

    \item By \Cref{lem:monotoneR}, we have that
    \[
        R_1 \subseteq R_2 \subseteq \dots \subseteq R_\ell.
    \]

    \item By \Cref{lem:default}, we have that for all $P,Q \in \QQ_i$ either
    \[
        PQ \setminus R_i \subseteq E(G)
        \quad
        \text{or}
        \quad
        PQ \setminus R_i \subseteq E(\overline{G}).
    \]
    \end{enumerate}
    Fix $r\in\N$, and let $k$ be the radius-$(5r+1)$ width of the given chain.
    For every $i\in\{2,\ldots,\ell\}$ and $v\in V(G)$, the definition of chain width gives
    \[
        |\Ball_{5r + 1}^{\PP_i}(v) / \PP_{i-1}| \leq k.
    \]
    Applying \Cref{lem:transversalRefinement} to
$\PP_{i-1}\preceq\PP_i$ and its transversal $T_{i-1}$ yields
    \begin{equation}\label{eq:ball-small}
         |\Ball_{5r}^{\PP_i}(v) / \QQ_{i-1} | = |\Ball_{5r}^{\PP_i}(v) / (\PP_{i-1} \wedge T_{i-1})| \leq k \cdot \pi_G(k).
    \end{equation}
    By \Cref{lem:ballContainment}, we have
    \begin{equation}\label{eq:ball-contained}
        \Ball_{r}^{R_i}(v)
        \subseteq
        \Ball_{5r}^{\PP_i}(v).
    \end{equation}
    Combining \eqref{eq:ball-small} and \eqref{eq:ball-contained} yields the desired bound on the radius-$r$ width of the merge sequence:
    \begin{equation*}
        |\Ball_{r}^{R_i}(v)/\QQ_{i-1}| \leq k \cdot \pi_G(k).
    \end{equation*}

    Finally, consider the graph $G$ equipped with two unary predicates marking $S_i$ and $T_i$.
    The relation $R_i$ is definable in this expansion by \Cref{lem:def-R} by a formula of quantifier rank $5$.
    The relation $\equiv_{\QQ_i}$ is
    defined easily using $S_i$ and $T_i$ by a formula of quantifier rank $2$.

    We now analyze the running time of a straightforward algorithm computing the resulting merge sequence. Put $n\coloneqq |V(G)|$.
    For each $i\in[\ell]$, keep every vertex of $S_i$ in a singleton part and group the remaining vertices by their adjacency vectors over $S_i$.
    This computes $\PP_i$ in $O(n(|S_i|+1))\le O(n^2)$ time.
    Scanning the resulting partition computes $T_i$ in $O(n)$ time.
    Similarly, compute the partition $\QQ_i$ in time $O(n(|T_i|+1))\le O(n^2)$.
    Thus computing and outputting $\PP_i$, $\QQ_i$, $S_i$, and $T_i$ takes $O(n^2)$ time for a fixed $i$.

    Compute $R_i$ by running breadth-first searches to depth five from every vertex in each of $G[X,Y]$ and $\overline G[X,Y]$, for all pairs $X,Y\in\PP_i$, allowing $X=Y$.
    The cost per pair is $O((|X|+|Y|)|X||Y|)$. We have:
    \[
        \sum_{X,Y\in\PP_i} (|X|+|Y|)|X||Y|
        =2\sum_{X,Y\in\PP_i} |X||Y|^2
        =2\sum_{X\in\PP_i} |X|\cdot \sum_{Y\in\PP_i}|Y|^2
        \le 2n(\sum_{Y\in\PP_i}|Y|)^2=2n^3,
    \]
     thus $R_i$ can be computed in time $O(n^3)$, and the whole merge sequence in time $O(\ell n^3)$.
\end{proof}
% lax end

\section{Construction sequences with locally definable partial graphs}\label{sec:construction-sequences}

For the purpose of model checking,
we need to work with construction sequences rather than merge sequences.
The construction of~\cite[Lemma~3.3]{mw-arxiv} shows how to convert a merge sequence to a construction sequence. However, a single step in a merge sequence is translated into a potentially long sequence of steps in a construction sequence. When aiming at defining each step of the construction sequence in the original graph $G$, we need to remember the result of the entire sequence of all the steps performed so far.
We cannot hope to obtain construction sequences in which each part is definable by a formula with fewer than $n/3$ parameters, even for the edgeless graph on $n$ vertices.
Therefore, in the following lemma, we only aim at proving that \emph{locally} -- in each ball of radius $r$ in the graph formed by the resolved vertex pairs -- the current set of resolved edges and each part of the current partition are definable in $G$ by formulas of bounded quantifier rank, involving a number of parameters which is bounded by the radius-$r$ merge-width of the input merge sequence. Recall the precise definition
of definability stated in \Cref{def:partial-definable}, and used in the statement below.
It is crucial that when considering merge sequences of bounded length, we
can retain the unary predicates describing every stage, and thus obtain formulas using only boundedly many unary predicates.

% lax begin Lax185574.MergeToConstruction
\begin{lemma}\label{cl:from-merge-to-construction}
    For every $n$-vertex graph $G$ and merge sequence $(\PP_1,R_1),\ldots,(\PP_\ell,R_\ell)$ of $G$, there is a construction sequence $G_1,\ldots,G_\lambda$ of $G$ such that the following properties hold for all $r\ge 1$.
    Let $w_r$ be the radius-$r$ width of the merge sequence.
    \begin{enumerate}
        \item The obtained construction sequence has radius-$r$ width at most $w_r$, and has length $\lambda\le O(nw_1)$.
        \item Suppose that for every $i\in[\ell]$, the relations $\equiv_{\PP_i}$ and $R_i$ are definable in some expansion of $G$ with~$2$ unary predicates, using two fixed first-order formulas of quantifier rank at most~$5$.

        Then, for every $t\in[\lambda-1]$ and  $v\in V(G)$, the induced partial graph
        $G_t[\Ball^{G_{t+1}}_r(v)]$ is definable in some expansion of $G$
        with $2\ell+1$ unary predicates, using formulas of quantifier rank~$7$ with $w_r$
        parameters.

        \item  The construction sequence can be
computed in time $O(\ell n^2)$, given an $n$-vertex graph $G$ and a merge sequence of length $\ell$.    \end{enumerate}
    \end{lemma}
% lax end

% lax begin Lax185574Proofs.MergeToConstruction.exists_constructionSequence_of_mergeSequence
\begin{proof}
We convert the merge sequence to a construction sequence by following the proof of~\cite[Lemma~3.3]{mw-arxiv}. Set $V\coloneqq V(G)$, $R_{\ell+1}\coloneqq {V\choose 2}$, and $\PP_{\ell+1}=\set{V}$.

Start with the singleton partition $\PP$ of $V$ and empty set $R$ of resolved vertex pairs.
The construction proceeds in \emph{stages} $j=1,\ldots,\ell$.
The invariant that is maintained during stage $j$ is that in each step, the current partition $\PP$ and current set $R$ of resolved vertex pairs satisfy:
\begin{align}\label{invariant-partitions}
    \PP_j&\preceq \PP\preceq \PP_{j+1}\qquad\\ R&\subseteq R_{j+1}\label{invariant-resolved}\\
    XY\setminus R\subseteq E(G)\quad&\text{ or }\quad XY\setminus R\subseteq E(\overline G)\quad\text{for all $X,Y\in \PP$}.
\end{align}

We start the process with $\PP \coloneqq \PP_1$ being the partition into singletons and $R \coloneqq \emptyset$.
Stage
$j$ proceeds as follows. First consider all unordered pairs of parts
$X,Y\in\PP_{j}$, allowing $X=Y$, for which
$XY\subseteq R_{j+1}$.
If there are unresolved vertex pairs in $XY$ (i.e.\ $XY\nsubseteq R$), resolve the pair $X,Y$
with the common sign -- positively if $XY\setminus R\subseteq E(G)$ and negatively if $XY\setminus R\subseteq E(\overline G)$ -- and update the set of resolved vertex pairs $R\coloneqq R\cup XY$.
Then repeatedly merge pairs $A,B$ of parts of the current partition for which $A\cup B$ is contained in one part of $\PP_{j+1}$, until the current partition $\PP$ becomes $\PP_{j+1}$; if $j=\ell$, terminate. This results in a construction sequence of $G$.

Fix $r\geq 1$.
We now upper bound the radius-$r$ width of the obtained construction sequence by the radius-$r$ width $w_r$ of the merge sequence.
Fix a stage $j$ in the merge sequence and a step
of the construction above in that phase, with current partition $\PP$ and current set of resolved pairs $R\subseteq {V\choose 2}$. By \eqref{invariant-resolved}, for every $v\in V$ we have:
\begin{align}\label{ball-inc}
    \Ball_r^{R}(v)&\subseteq \Ball_r^{R_{j+1}}(v).
\end{align}
As $\Ball_r^{R_{j+1}}(v)$ intersects at most $w_r$ parts of $\PP_j$ and $\PP$ is coarser than $\PP_j$ by \eqref{invariant-partitions},
also $\Ball_r^{R}(v)$ intersects at most $w_r$ parts of $\PP$, proving the required bound.

The obtained construction sequence is \emph{effective}: each resolve operation introduces at least one new resolved vertex pair. By~\cite[Lemma~3.4]{mw-arxiv} (and the comment before it), the length $\lambda$ of the construction sequence is at most $n(2w_1+1) \leq O(nw_1)$.

A straightforward implementation processes each stage in time $O(n^2)$.
Namely, for each pair $X,Y$ of parts of $\PP_j$,
test whether $XY\subseteq R_{j+1}$ by scanning
all vertex pairs $uv\in XY$, and record the sign with which $X,Y$ should be resolved. As this processes each pair $uv\in {V\choose 2}$ once, this takes time $O(n^2)$ in total. Similarly, the merge phase can be implemented in time $O(n)$.
Altogether, the  construction sequence is computable in time $O(\ell n^2)$.
This proves the first and last conditions of the lemma.

\medskip

We now prove the second condition of the lemma, regarding definability.
Fix a step $G_t\to G_{t+1}$ of the obtained construction sequence, and let $j$ be the stage
in which this step is performed.
Let $\PP$ and $R$ denote the partition and the set of resolved vertex pairs in $G_t$, and let $R_{+}$ denote the set of resolved vertex pairs in $G_{t+1}$; then $R\subseteq R_{+}\subseteq R_{j+1}$.

  Let $\phi_\equiv^i(x,y)$ and $\phi_R^i(x,y)$ be formulas defining
$\equiv_{\PP_i}$ and $R_i$, respectively, as given by the assumption.
For $1\le j<\ell$, the set $$\bigcup\setof{XY}{X,Y\in\PP_j, XY\subseteq R_{j+1}}$$ is defined
by the formula
\[
 \gamma_j(x,y) \,\coloneqq\, x\ne y\ \land\forall a\forall b\,\bigl(
       (\phi_\equiv^{j}(x,a)\land\phi_\equiv^{j}(y,b)\land a\ne b)\,\to\,\phi_R^{j+1}(a,b)\bigr),
\]
which has quantifier rank at
most~$7$.  Consequently, at the start of stage~$j$, the resolved
relation, which we denote $R_j'$\,,
is defined by
\[
 \Gamma_{j}(x,y)\coloneqq \bigvee_{i=1}^{j-1}\gamma_i(x,y).
\]
This formula uses $2j$ unary predicates and has quantifier rank at most $7$.

Choose
one parameter $p_X\in X$ for every part $X\in \PP_j$.
The part $X$ is then defined by $\phi_\equiv^{j}(x,p_X)$, which
uses $2$ unary predicates, one parameter, and has quantifier rank at most $5$.

\medskip
Fix a vertex $v\in V$.
Denote
\[
 B\coloneqq \Ball^{R_{+}}_r(v)=\Ball^{G_{t+1}}_r(v)\subseteq\Ball^{R_{j+1}}_r(v).
\]
Mark the set $B$ with a unary predicate $U_B$.

We will show how to define in $G$ the relation $R\cap {B\choose 2}$ and the sets $X\cap B$, for each $X\in \PP$,
using Boolean combinations of the formulas defined above,
using the $2j+1$ distinguished unary predicates, and at most $w_r$ parameters.
Once $R\cap {B\choose 2}$ is defined, the set of  edges and non-edges of $G_t$, restricted to $B\choose 2$, can be recovered by further intersecting with $E(G)$ and $E(\overline G)$, respectively. This will therefore define the induced substructure $G_t[B]$, as required in the second condition of the lemma.

\medskip
We first focus on defining $R\cap {B\choose 2}$.
Note that $R$
is the union of $R_j'$
and all already processed rectangles $XY$ with $X,Y\in\PP_j, XY\subseteq R_{j+1}$.
Therefore, $R$ is defined by the disjunction
of $\Gamma_j(x,y)$ and disjuncts of the form
\[
 (x\neq y)\land \left((\phi_\equiv^{j}(x,p_X)\land\phi_\equiv^{j}(y,p_Y))
 \lor
 (\phi_\equiv^{j}(x,p_Y)\land\phi_\equiv^{j}(y,p_X))\right).
\]
To define the intersection of $R$ with $B\choose 2$,
add a conjunct $U_B(x)\land U_B(y)$,
and keep only those disjuncts which involve parts $X,Y$ that intersect $B$. As there are at most $w_r$ many such parts, this results in a formula
 using $w_r$ parameters.

Each current part $X\in\PP$ is a union of parts of $\PP_j$, so its intersection with $B$ is defined by
the formula
$$U_B(x)\land \bigvee_{Y}\phi_\equiv^{j}(x,p_Y)$$ where $Y$ ranges over all parts $Y\in\PP_j$ with $Y\cap B\neq\emptyset$ and $Y\subseteq X$, with empty disjunction interpreted as \emph{false}.
The resulting formula uses only $w_r$ parameters.
This completes the proof of the lemma.
\end{proof}
% lax end

We now prove \Cref{lem:better-definable-mw-from-sets}.
\lembetterdefinablemwfromsets*

% lax begin Lax185574Proofs.ChainConstructionSequences.exists_constructionSequence_of_chain
\begin{proof}
    The empty graph is immediate, so assume $G$ is nonempty.
    Apply \Cref{lem:better-mw} to the given chain to obtain a merge sequence with $\ell$ stages, each definable by two unary predicates and fixed formulas of quantifier rank at most $5$.
    Fix $r\ge 1$ and let $k_r$ be the radius-$(5r+1)$ width of the chain.
    Then the radius-$r$ width $w_r$ of this merge sequence is at most $k_r\cdot \pi_G(k_r)$. Apply \Cref{cl:from-merge-to-construction} to obtain a construction sequence of length $\lambda\le O(nw_1)$ and radius-$r$ width at most $w_r$, with local definitions using at most $w_r$ parameters, $2\ell+1$ unary predicates, and formulas of quantifier rank at most $7$.
    Both algorithms take $O(\ell n^3)$ time and use only $G$ and the given chain.
\end{proof}
% lax end
